% Session 5 -- engineering a Kitaev chain out of a quantum spin Hall edge.
% Placement: sec. "Engineering a topological superconductor from a quantum spin Hall insulator",
% after the proximity-coupling paragraph.
%
% The two Majoranas sit at the ends of the proximitized segment, i.e. exactly at the two
% magnet/superconductor interfaces, per the text.
\begin{aqmfigure}{fig:s05:qsh}{Building the Kitaev chain out of parts that exist. The edge of a
quantum spin Hall insulator is helical: on one and the same edge, spin-up electrons run one way and
spin-down electrons the other. Depositing a ferromagnet on a stretch of that edge breaks
time-reversal symmetry locally and destroys the Kramers protection there, gapping the edge
underneath it; two ferromagnets leave one short helical segment surviving between them. That
segment offers only one propagation direction per spin, so counted as channels it is a
\emph{spinless} one-dimensional wire -- precisely the ingredient (5.19) demands and the one nature
does not hand us directly. Proximity coupling it to an ordinary $s$-wave superconductor supplies
the pairing, and the result is a one-dimensional $p$-wave superconductor carrying a Majorana zero
mode at each end of the segment, that is, at the two ferromagnet--superconductor interfaces. Note
what was \emph{not} required: no unconventional superconductor appears anywhere in the recipe, only
an $s$-wave one, a magnet, and a topological insulator.}

\begin{tikzpicture}[x=1cm, y=1cm]
\tikzset{
  gline/.style={aqm, draw=aqmmuted, line width=0.6pt, dash pattern=on 2pt off 1.5pt},
  chan/.style={aqm, line width=1.1pt, -{Stealth[length=5pt,width=4pt]}},
  majdot/.style={aqm, circle, draw=aqmbroken, fill=aqmfillb, line width=1.1pt,
                 minimum size=10pt, inner sep=0pt}
}

% ============================ the deposited layer ============================
% ferromagnet | superconductor | ferromagnet, all sitting on top of the same edge
\foreach \xa/\xb in {0/4.30, 8.70/13.0} {
  \fill[white] (\xa,2.45) rectangle (\xb,3.35);
  \path[pattern=north east lines, pattern color=aqmmuted] (\xa,2.45) rectangle (\xb,3.35);
  \draw[aqm, draw=aqmmuted] (\xa,2.45) rectangle (\xb,3.35);
  \node[aqmlab, anchor=center, fill=white, inner sep=2.5pt] at ({(\xa+\xb)/2},2.90)
    {ferromagnet};
}
\fill[aqmfill] (4.30,2.45) rectangle (8.70,3.35);
\draw[aqm, draw=aqmaccent] (4.30,2.45) rectangle (8.70,3.35);
\node[aqmlab, anchor=center] at (6.50,2.90) {$s$-wave superconductor};

% ============================ the interfaces =================================
\draw[gline] (4.30,0.36) -- (4.30,2.43);
\draw[gline] (8.70,0.36) -- (8.70,2.43);

% ============================ what each block does ===========================
\node[aqmlabsmall, anchor=north, align=center, text=aqmmuted]
  at (2.15,2.28) {time reversal broken:\\[-2pt]helical modes gapped};
\node[aqmlabsmall, anchor=north, align=center, text=aqmmuted]
  at (10.85,2.28) {time reversal broken:\\[-2pt]helical modes gapped};
\node[aqmlabsmall, anchor=north, align=center]
  at (6.50,2.28) {one direction per spin:\\[-1pt]\textbf{effectively spinless}};

% ============================ the helical edge ===============================
% the surviving segment
\draw[chan, draw=aqmbroken] (4.55,1.15) -- (8.45,1.15);
\node[aqmlabsmall, fill=white, inner sep=1.5pt, text=aqmbroken]
  at (6.50,1.15) {$\uparrow$ moves right};
\draw[chan, draw=aqmbroken] (8.45,0.75) -- (4.55,0.75);
\node[aqmlabsmall, fill=white, inner sep=1.5pt, text=aqmbroken]
  at (6.50,0.75) {$\downarrow$ moves left};
% the gapped stretches
\foreach \xa/\xb in {0.20/4.05, 8.95/12.80} {
  \draw[gline] (\xa,1.15) -- (\xb,1.15);
  \draw[gline] (\xa,0.75) -- (\xb,0.75);
}

% ============================ the slab =======================================
\fill[aqmfill] (0,-1.30) rectangle (13.0,0.36);
\draw[aqm, draw=aqmaccent] (0,-1.30) rectangle (13.0,0.36);
\node[aqmlab, anchor=center] at (6.50,-0.47)
  {quantum spin Hall insulator, gapped bulk \; (e.g.\ WTe$_2$)};

% ============================ the Majoranas ==================================
\node[majdot] at (4.30,0.95) {};
\node[majdot] at (8.70,0.95) {};
\node[aqmlabsmall, anchor=east, text=aqmbroken] at (4.08,0.95) {$\gamma_1$};
\node[aqmlabsmall, anchor=west, text=aqmbroken] at (8.92,0.95) {$\gamma_2$};

% ============================ the read-off ===================================
\node[aqmboxbroken, anchor=north, align=center] at (6.50,-1.60)
  {an effectively spinless 1D segment $+$ proximity pairing $=$ the Kitaev chain (5.19),\\[2pt]
   with its two unpaired end Majoranas now pinned\\[-1pt]
   at the ferromagnet--superconductor interfaces};

\end{tikzpicture}

\end{aqmfigure}
